% GENERATED FILE. Edit canonical lessons, then run npm run generate:books.
% canonical-source-hash: fnv1a64:f7fc1f8520a38682
% canonical-lessons: GE-C473-kichern, GE-C473-grinsen, GE-C473-schmunzeln, GE-C473-albern, GE-C473-kitzeln, GE-R473-review-466-473

\chapter{Kichern und grinsen}
\label{ch:ge-kichern-und-grinsen}
\hlchaptermodality{473}

\begin{chapteropening}
\textbf{By the end of this chapter:} \emph{I can describe different ways of laughing and smiling.}

Complete the last lesson of chapter 473: I can describe different ways of laughing and smiling.
\end{chapteropening}

\section[kichern]{kichern — to giggle}

\label{lesson:GE-C473-kichern}



\begin{quote}
Before the new one: say the German for homesickness, then the German for to miss.
\end{quote}

\subsection*{You'll want to know: kichern}
\textbf{kichern} — \textquotedblleft{}to giggle\textquotedblright{}.

\begin{quote}
Die Mädchen kicherten in der letzten Reihe.
\end{quote}

\begin{quote}
\emph{The girls were giggling in the back row.}
\end{quote}

A softer, higher laugh than \emph{lachen}.

\begin{cousinweb}
\emph{kichern} is usually explained as an imitative word, copying the sound of a high, held-back laugh; EtymDB lists it without a source, so its origin is not settled.
\end{cousinweb}

\begin{grammarlens}[title={What you've built}]
One more word for telling, explaining and describing at B1.
\end{grammarlens}

\subsection*{Your turn}
\begin{itemize}
\raggedright
  \item \emph{Say it:} \emph{kichern}
  \item \emph{Say it:} \emph{kichern}, once more
  \item \emph{Say it:} a sentence of your own with \emph{kichern}
  \item \emph{Recall:} say the German for homesickness, then the German for to miss, then say \emph{kichern} again
\end{itemize}

\begin{tcolorbox}[breakable,colback=teal!4,colframe=teal!35!black,title={Before you move on}]
What does kichern mean? (\textquotedblleft{}To giggle\textquotedblright{}.) Say it once more.
\end{tcolorbox}
\section[grinsen]{grinsen — to grin}

\label{lesson:GE-C473-grinsen}



\begin{quote}
Before the new one: say the German for to miss, then the German for to giggle.
\end{quote}

\subsection*{You'll want to know: grinsen}
\textbf{grinsen} — \textquotedblleft{}to grin\textquotedblright{}.

\begin{quote}
Er grinste, als er seine gute Note sah.
\end{quote}

\begin{quote}
\emph{He grinned when he saw his good mark.}
\end{quote}

Often a little cheeky: \emph{frech grinsen}.

\begin{cousinweb}
\emph{grinsen} is usually traced to Middle High German \emph{grinsen}, beside \emph{grīnen} \textquotedblleft{}to bare the teeth\textquotedblright{}; its older origin is not settled, and EtymDB has no entry.
\end{cousinweb}

\begin{grammarlens}[title={What you've built}]
One more word for telling, explaining and describing at B1.
\end{grammarlens}

\subsection*{Your turn}
\begin{itemize}
\raggedright
  \item \emph{Say it:} \emph{grinsen}
  \item \emph{Say it:} \emph{grinsen}, once more
  \item \emph{Say it:} a sentence of your own with \emph{grinsen}
  \item \emph{Recall:} say the German for to miss, then the German for to giggle, then say \emph{grinsen} again
\end{itemize}

\begin{tcolorbox}[breakable,colback=teal!4,colframe=teal!35!black,title={Before you move on}]
What does grinsen mean? (\textquotedblleft{}To grin\textquotedblright{}.) Say it once more.
\end{tcolorbox}
\section[schmunzeln]{schmunzeln — to smile quietly, to smirk}

\label{lesson:GE-C473-schmunzeln}



\begin{quote}
Before the new one: say the German for to giggle, then the German for to grin.
\end{quote}

\subsection*{You'll want to know: schmunzeln}
\textbf{schmunzeln} — \textquotedblleft{}to smile quietly, to smirk\textquotedblright{}.

\begin{quote}
Der Lehrer schmunzelte über die Antwort.
\end{quote}

\begin{quote}
\emph{The teacher smiled to himself at the answer.}
\end{quote}

A quiet, amused smile with the lips closed.

\begin{cousinweb}
\emph{schmunzeln} is traced in EtymDB to Middle High German \emph{smutzen} \textquotedblleft{}to smile\textquotedblright{}; how the modern form arose is not settled.
\end{cousinweb}

\begin{grammarlens}[title={What you've built}]
One more word for telling, explaining and describing at B1.
\end{grammarlens}

\subsection*{Your turn}
\begin{itemize}
\raggedright
  \item \emph{Say it:} \emph{schmunzeln}
  \item \emph{Say it:} \emph{schmunzeln}, once more
  \item \emph{Say it:} a sentence of your own with \emph{schmunzeln}
  \item \emph{Recall:} say the German for to giggle, then the German for to grin, then say \emph{schmunzeln} again
\end{itemize}

\begin{tcolorbox}[breakable,colback=teal!4,colframe=teal!35!black,title={Before you move on}]
What does schmunzeln mean? (\textquotedblleft{}To smile quietly, to smirk\textquotedblright{}.) Say it once more.
\end{tcolorbox}
\section[albern]{albern — silly}

\label{lesson:GE-C473-albern}



\begin{quote}
Before the new one: say the German for to grin, then the German for to smile quietly.
\end{quote}

\subsection*{You'll want to know: albern}
\textbf{albern} — \textquotedblleft{}silly\textquotedblright{}.

\begin{quote}
Die Kinder waren müde und deshalb albern.
\end{quote}

\begin{quote}
\emph{The children were tired and therefore silly.}
\end{quote}

In speech also \emph{herumalbern} — to fool around.

\begin{cousinweb}
\emph{albern} is inherited from Middle High German \emph{alwære} and Old High German \emph{alawāri} \textquotedblleft{}friendly, good-natured\textquotedblright{} (EtymDB); the meaning sank from kind to simple to silly.
\end{cousinweb}

\begin{grammarlens}[title={What you've built}]
One more word for telling, explaining and describing at B1.
\end{grammarlens}

\subsection*{Your turn}
\begin{itemize}
\raggedright
  \item \emph{Say it:} \emph{albern}
  \item \emph{Say it:} \emph{albern}, once more
  \item \emph{Say it:} a sentence of your own with \emph{albern}
  \item \emph{Recall:} say the German for to grin, then the German for to smile quietly, then say \emph{albern} again
\end{itemize}

\begin{tcolorbox}[breakable,colback=teal!4,colframe=teal!35!black,title={Before you move on}]
What does albern mean? (\textquotedblleft{}Silly\textquotedblright{}.) Say it once more.
\end{tcolorbox}
\section[kitzeln]{kitzeln — to tickle}

\label{lesson:GE-C473-kitzeln}



\begin{quote}
Before the new one: say the German for to smile quietly, then the German for silly.
\end{quote}

\subsection*{You'll want to know: kitzeln}
\textbf{kitzeln} — \textquotedblleft{}to tickle\textquotedblright{}.

\begin{quote}
Der Vater kitzelt seine Tochter, und sie lacht laut.
\end{quote}

\begin{quote}
\emph{The father tickles his daughter, and she laughs out loud.}
\end{quote}

\emph{Das kitzelt!} — That tickles!

\begin{cousinweb}
\emph{kitzeln} is inherited from Old High German \emph{kizzilōn} \textquotedblleft{}to tickle\textquotedblright{} (EtymDB).
\end{cousinweb}

\begin{grammarlens}[title={What you've built}]
One more word for telling, explaining and describing at B1.
\end{grammarlens}

\subsection*{Your turn}
\begin{itemize}
\raggedright
  \item \emph{Say it:} \emph{kitzeln}
  \item \emph{Say it:} \emph{kitzeln}, once more
  \item \emph{Say it:} a sentence of your own with \emph{kitzeln}
  \item \emph{Recall:} say the German for to smile quietly, then the German for silly, then say \emph{kitzeln} again
\end{itemize}

\begin{tcolorbox}[breakable,colback=teal!4,colframe=teal!35!black,title={Before you move on}]
What does kitzeln mean? (\textquotedblleft{}To tickle\textquotedblright{}.) Say it once more.
\end{tcolorbox}
\section[(dialogue)]{Gefühle und Gesundheit}

\label{lesson:GE-R473-review-466-473}



\begin{quote}
Say the German for jealousy and envy.
\end{quote}

\subsection*{Your turn}
Say these aloud:

\begin{itemize}
\raggedright
  \item jealousy, envy, envious, to envy, to not begrudge
  \item shame, to be ashamed, embarrassed, embarrassment, to blush
  \item rage, anger, to annoy, to get worked up, to rage
  \item to disappoint, disappointment, relief, to breathe a sigh of relief, to sigh
  \item patience, calm, to calm, excitement, to fidget
  \item tear, to sob, comfort, to comfort, pity
  \item to long, longing, loneliness, homesickness, to miss
  \item to giggle, to grin, to smile quietly, silly, to tickle
\end{itemize}

\begin{tcolorbox}[breakable,colback=teal!4,colframe=teal!35!black,title={Before you move on}]
Jealousy? (\textbf{die Eifersucht}.) Envy? (\textbf{der Neid}.) Envious? (\textbf{neidisch}.) To envy? (\textbf{beneiden}.) To not begrudge? (\textbf{gönnen}.) Shame? (\textbf{die Scham}.) To be ashamed? (\textbf{schämen}.) Embarrassed? (\textbf{verlegen}.) Embarrassment? (\textbf{die Verlegenheit}.) To blush? (\textbf{erröten}.) Rage? (\textbf{die Wut}.) Anger? (\textbf{der Zorn}.) To annoy? (\textbf{ärgern}.) To get worked up? (\textbf{aufregen}.) To rage? (\textbf{toben}.) To disappoint? (\textbf{enttäuschen}.) Disappointment? (\textbf{die Enttäuschung}.) Relief? (\textbf{die Erleichterung}.) To breathe a sigh of relief? (\textbf{aufatmen}.) To sigh? (\textbf{seufzen}.) Patience? (\textbf{die Geduld}.) Calm? (\textbf{gelassen}.) To calm? (\textbf{beruhigen}.) Excitement? (\textbf{die Aufregung}.) To fidget? (\textbf{zappeln}.) Tear? (\textbf{die Träne}.) To sob? (\textbf{schluchzen}.) Comfort? (\textbf{der Trost}.) To comfort? (\textbf{trösten}.) Pity? (\textbf{das Mitleid}.) To long? (\textbf{sehnen}.) Longing? (\textbf{die Sehnsucht}.) Loneliness? (\textbf{die Einsamkeit}.) Homesickness? (\textbf{das Heimweh}.) To miss? (\textbf{vermissen}.) To giggle? (\textbf{kichern}.) To grin? (\textbf{grinsen}.) To smile quietly? (\textbf{schmunzeln}.) Silly? (\textbf{albern}.) To tickle? (\textbf{kitzeln}.)
\end{tcolorbox}
